\documentclass[11pt]{article}
\usepackage[margin=1in]{geometry}
\usepackage{amsmath,amssymb,amsthm,mathtools}
\usepackage[T1]{fontenc}
\usepackage[colorlinks=true,allcolors=blue]{hyperref}
\newtheorem{theorem}{Theorem}
\newtheorem{lemma}[theorem]{Lemma}
\newcommand{\Herm}{\operatorname{Herm}}
\newcommand{\conv}{\operatorname{conv}}
\newcommand{\tr}{\operatorname{tr}}
\newcommand{\Sep}{\operatorname{Sep}}
\title{An application of separable-state nonrepresentability\ to Grassmann orbitopes and calibrations}
\author{Alec Kriebel\thanks{Independent researcher. ORCID:
\href{https://orcid.org/0009-0001-9320-500X}{0009-0001-9320-500X}.}}
\date{12 October 2026; version 1.0; unrefereed preprint}
\begin{document}
\maketitle
\begin{abstract}
We apply Fawzi's established obstruction for normalized complex separable
states to the oriented real Grassmann orbitope $C_{4,12}$. Two affine
hyperplanes select a section real-linearly isomorphic to $\Sep(3,3)$,
so this mass ball has no finite projected semidefinite lift. A compact
first-moment argument then disproves the universal degree-one supporting
sum-of-squares assertion formulated by Sanyal, in joint work with Rostalski.
It also gives a negative answer in general to Harvey--Lawson's Question 6.5:
some constant comass-one four-form on $\mathbb R^{12}$ has no pointwise or
local norm completion. The failing unnormalized functional may have rational
coefficients, but no explicit coefficient list or exact comass is supplied.
The reduction combines classical K\"ahler faces, established fermionic
projector bodies and the prior tensor-to-wedge lemma with Fawzi's theorem;
worldwide priority of this assembly is not certified.
\end{abstract}

\section{The original question and the result}
Let $X_{k,n}$ consist of all decomposable unit real $k$-vectors in
$\Lambda^k\mathbb R^n$, including both orientations, and let
$C_{k,n}=\conv(X_{k,n})$. Write $I(X_{k,n})$ for the full real vanishing
ideal. Sanyal's report~\cite[pp.708--709]{sanyal} asks whether, for every
real linear functional $\ell$ and $\delta=\max_{X_{k,n}}\ell$, one has
\begin{equation}\label{eq:sos}
 \delta-\ell=\sum_j h_j^2\pmod{I(X_{k,n})},\qquad \deg h_j\leq1.
\end{equation}
Thus the $h_j$ may be affine, including constants. That source credits
joint computational work with Philipp Rostalski and the equivalence with
Harvey--Lawson's Question 6.5~\cite[p.68]{hl}. This connection and the
first-theta-body formulation already appear in Rostalski's 2010 slides and
Gouveia--Thomas's chapter~\cite{rostalski,chapter}.

Here a finite semidefinite lift means a linear image of a spectrahedron of
any finite size, with arbitrary affine equalities and auxiliary variables.
It is stronger to exclude such lifts than to exclude an unprojected LMI or
one specified hierarchy. We use the following established theorem: the
normalized complex body
\[
 \Sep(3,3)=\conv\{xx^*\otimes yy^*:x,y\in\mathbb C^3,
                      \|x\|=\|y\|=1\}
\]
has no finite semidefinite lift. This is the $3\times3$ case of Fawzi's
Theorem 1~\cite[pp.1320--1321]{fawzi}, public in its 2019 preprint. His
definition allows projected Hermitian LMIs and therefore also excludes
real-symmetric lifts. We do not reprove that external theorem.

\begin{theorem}\label{thm:main}
The body $C_{4,12}$ has no finite semidefinite lift. Consequently
\eqref{eq:sos} fails for some linear functional on $\Lambda^4\mathbb R^{12}$,
which may be chosen to have rational coefficients. Some constant closed
comass-one four-form on $\mathbb R^{12}$ has no norm completion by finitely
many four-forms, even at a single point.
\end{theorem}

This refutes the original universal degree-one assertion and the general
affirmative answer to Question 6.5. The conclusion is existential. We claim
no explicit failing form, exact value of its maximum, rational normalized
calibration, least dimension or lower-dimensional classification.

\section{The K\"ahler face as a projector body}
Identify $V=\mathbb R^{12}$ with the underlying real space of
$E=\mathbb C^6$. Fix an orthonormal real basis $e_1,f_1,\ldots,e_6,f_6$
with $Je_j=f_j$, $Jf_j=-e_j$. Set
\[
 \omega=\sum_{j=1}^6 e_j^*\wedge f_j^*,\qquad \Phi=\omega^2/2.
\]
The classical Wirtinger inequality gives $\Phi(p)\leq1$ for $p\in X_{4,12}$,
with equality precisely for complex two-planes with their complex
orientation~\cite{hl,barvinok}. For clarity, the compression of $J$ to a
real four-plane has two skew singular values of magnitude at most one.
The Pfaffian evaluates $\Phi$ as their signed product. Equality forces
the compression to preserve every vector's norm, hence $J$ preserves the
plane; the positive Pfaffian selects its complex orientation. Conversely,
a complex-orthonormal pair $a,b$ gives the real frame $a,Ja,b,Jb$ of value
one. Thus
\[
 F=C_{4,12}\cap\{\Phi=1\}
\]
is exactly the convex hull of these positively oriented complex planes:
every positive-weight summand in a convex expression of value one must
itself attain one.

Let $W=\Lambda^2_{\mathbb C}E$, of complex dimension 15, with its induced
Hermitian metric. In the complexification of $V$, put
$u_j=(e_j-if_j)/\sqrt2$, $u_I=u_i\wedge u_j$ for $I=(i,j)$, $i<j$.
Define a real-linear map
\begin{equation}\label{eq:map}
 A:\Herm(W)\longrightarrow\Lambda^4_{\mathbb R}V,
 \qquad A(H)=\sum_{I,J}H_{IJ}u_I\wedge\overline{u_J}.
\end{equation}
The result is real: conjugation transposes $H$, while swapping exterior
factors of degree two has sign $+1$. The 225 vectors
$u_I\wedge\overline{u_J}$ form a complex basis of type $(2,2)$, so $A$ is
injective on the 225-dimensional real Hermitian space. With
$w_j=e_j-if_j$, the same formula is
$A(H)=\tfrac14\sum H_{IJ}w_I\wedge\overline{w_J}$; its real coefficients
on the standard Hermitian basis are rational.

For $a\in E$, write $\widehat a=\sum a_j u_j$. Its underlying real vector
satisfies
\[
 a=(\widehat a+\overline{\widehat a})/\sqrt2,\quad
 Ja=i(\widehat a-\overline{\widehat a})/\sqrt2,\quad
 a\wedge Ja=-i\widehat a\wedge\overline{\widehat a}.
\]
For $z=a\wedge b\in W$ this gives
\begin{equation}\label{eq:projector}
 A(zz^*)=\widehat a\wedge\widehat b\wedge
           \overline{\widehat a}\wedge\overline{\widehat b}
        =a\wedge Ja\wedge b\wedge Jb.
\end{equation}
The two factors $-i$ give $-1$, and moving $\overline{\widehat a}$ past
$\widehat b$ cancels that sign. A decomposable unit $z$ has a
complex-orthonormal factorization, absorbing its unit phase in a factor.
Conversely, every generating plane of $F$ has this form. Therefore
\begin{equation}\label{eq:fermion}
 F=A(K),\qquad
 K=\conv\{zz^*:z\in W\text{ decomposable},\ \|z\|=1\}.
\end{equation}
The body $K$ is the established mixed fermionic projector hull
\cite[Definition 4.2 and Theorem 5.1]{fermion}, rather than a new notion of
separability. Equation~\eqref{eq:fermion} is a real-linear isomorphism onto
the image. Evaluation on the Hermitian basis gives
$\Phi(A(H))=\tr H$: the coefficient on
$e_i\wedge f_i\wedge e_j\wedge f_j$ is $H_{(i,j),(i,j)}$.

\section{The normalized complex separable-state section}
Split $E=U\oplus Z$ orthogonally, using its first and last three complex
coordinates. Then
\[
 W=\Lambda^2U\oplus B\oplus\Lambda^2Z,
 \qquad B=U\wedge Z,\quad D=\Lambda^2U\oplus\Lambda^2Z.
\]
Let $P_D$ be the orthogonal projector onto $D$, and define
$K_0=\{H\in K:\tr(P_DH)=0\}$. If
$H=\sum_t\lambda_t z_tz_t^*$, with $\lambda_t>0$ and $\sum_t\lambda_t=1$,
then
\[
 0=\tr(P_DH)=\sum_t\lambda_t\|P_Dz_t\|^2.
\]
Each $z_t$ therefore lies in $B$. Positivity excludes cancellation and
proves the full generator statement, rather than only matrix support.

The unitary map $T:U\otimes_{\mathbb C}Z\to B$ defined by
$T(x\otimes y)=x\wedge y$ takes rank-one tensors exactly to nonzero
decomposable bivectors in $B$. This is a specialization of the prior
tensor-to-wedge theorem of Vrana--L\'evay~\cite[\S3.1, equation (25)]{vrana}.
Its elementary proof here is short. If $a\wedge b\in B$ is nonzero, the
$U$ projections of $a,b$ are dependent, as are their $Z$ projections,
because both pure exterior components vanish. Neither projection span
is zero, or the wedge would vanish. Taking generators $x\in U,y\in Z$,
write $a=\alpha x+\gamma y$, $b=\beta x+\delta y$. Their wedge is
$(\alpha\delta-\beta\gamma)x\wedge y$, with nonzero coefficient.
The converse is immediate. Orthogonality gives
$\|x\wedge y\|=\|x\|\|y\|$; for unit $z$ both factors can be rescaled to
unit norm. Under $T$ the projector becomes $xx^*\otimes yy^*$. Hence
\begin{equation}\label{eq:section}
 K_0\cong\Sep(3,3).
\end{equation}

This is an affine section in the original real exterior space. Put
\[
 \omega_U=\sum_{j=1}^3e_j^*\wedge f_j^*,\quad
 \omega_Z=\sum_{j=4}^6e_j^*\wedge f_j^*,\quad
 Q=(\omega_U^2+\omega_Z^2)/2.
\]
The same coefficient calculation gives $Q(A(H))=\tr(P_DH)$, so
\begin{equation}\label{eq:ambientsection}
 C_{4,12}\cap\{\Phi=1,\ Q=0\}=A(K_0)\cong\Sep(3,3).
\end{equation}
We need $Q\geq0$ only on $F$; it need not be nonnegative on the whole
mass ball. The equations are simultaneous, and membership in the mass
ball is essential. They alone do not define the separable set.

Affine sections and linear maps preserve finite semidefinite lifts:
additional affine equations on a projected point can be imposed on the
lift variables. A left inverse of $A$ on its image extends linearly to
the ambient space. Thus a lift of $C_{4,12}$ would give a lift of
$\Sep(3,3)$ through~\eqref{eq:ambientsection}, contradicting Fawzi.
This proves the first assertion of Theorem~\ref{thm:main}.

\section{Supporting affine squares and calibration completion}
We give the required compact moment argument explicitly. It specializes
the established theta-body and lift framework~\cite{theta,lifts,survey}.
Let $X\subset S^{N-1}\subset\mathbb R^N$ be nonempty and compact, and set
$J_2=I(X)\cap\mathbb R[p]_{\leq2}$. Choose a finite basis of this vector
space. For a symmetric $(N+1)$-square matrix $M$, define
\[
 L_M(1)=M_{00},\quad L_M(p_i)=M_{0i},\quad L_M(p_ip_j)=M_{ij}.
\]
Consider the spectrahedron and its projection
\[
 S_X=\{M\succeq0:M_{00}=1,\ L_M(g)=0\ (g\in J_2)\},\qquad
 T_X=\{(M_{01},\ldots,M_{0N}):M\in S_X\}.
\]
Because $1-\sum p_i^2\in J_2$, every feasible matrix has trace two.
Positive semidefiniteness bounds all its entries; thus $S_X$ and $T_X$
are compact. Rank-one matrices $(1,p)(1,p)^T$ for $p\in X$ imply
$\conv(X)\subseteq T_X$.

If every supporting affine function were an affine SOS on $X$, then
for $h_j(p)=a_j^T(1,p)$ the discrepancy
$\delta-\ell-\sum_jh_j^2$ would lie in $J_2$. Every feasible $M$, projecting
to $x$, would satisfy
\[
 \delta-\ell(x)=\sum_j a_j^TMa_j\geq0.
\]
Intersecting all supporting halfspaces would give $T_X=\conv(X)$,
a finite lift. There is no appeal to SDP duality, Slater conditions,
or closure of a nonclosed shadow.

Apply this to $X=X_{4,12}\subset S^{494}$, with $N=495$. Since the mass
ball has no finite lift, $T_X$ strictly contains it. Choose
$x\in T_X\setminus C_{4,12}$ and strictly separate it by a coefficient
vector $c$:
\[
 c\cdot x>\max_{p\in X}c\cdot p.
\]
The positive gap persists under small perturbations, since $x$ and $X$
are bounded. Choose a rational vector $c'$ in that neighborhood and
put $\ell(p)=c'\cdot p$, $\delta=\max_X\ell$. A lifting matrix for $x$
makes $\delta-\ell(x)<0$, which the preceding SOS calculation forbids.
Thus a rational unnormalized failing functional exists. Signed coordinate
wedges belong to $X$ and span the exterior space, so nonzero $\ell$ has
$\delta>0$. The constant form $\varphi=\ell/\delta$ is closed and has
comass one, including absolute-value normalization because $X=-X$.

If linear forms $\psi_j$ on the exterior space satisfied the norm
completion in Question 6.5, then on unit decomposable vectors
\[
 \varphi(p)^2+\sum_j\psi_j(p)^2=1
 \quad\Longrightarrow\quad
 1-\varphi(p)=\tfrac12(1-\varphi(p))^2+
                         \tfrac12\sum_j\psi_j(p)^2.
\]
This affine SOS contradicts the chosen functional. A local completion by
differential four-forms would give such a pointwise identity at each point,
so none can exist. This proves the remaining assertions.

For clarity on the prior equivalence, the converse at a point is also
elementary. If $1-\varphi(p)=\|a+Bp\|^2$ on $X$, evaluation at $\pm p$
and spanning give $B^Ta=-\varphi/2$ and
$\|a\|^2+\|Bp\|^2=1$. At a maximizer $p_0$, $Bp_0=-a$; evaluation at
$-p_0$ gives $\|a\|^2=1/2$. Projecting onto $a^\perp$ yields
$\|\sqrt2 P_{a^\perp}Bp\|^2=1-\varphi(p)^2$. Homogeneity extends this
to all simple vectors. No assertion about choosing smooth local completions
from arbitrary pointwise solutions is required.

\section{Related results, verification and limits}
Nonrepresentability for quadratic Grassmann cones is prior:
Bettiol--Kummer--Mendes~\cite[Corollary 4.2 and Remark 4.3]{curvature} treat
nonnegative homogeneous quadratics in
$\operatorname{Sym}^2(\Lambda^k\mathbb R^n)$. Blekherman--Smith--Velasco
\cite[Theorem 1.1]{bsv} characterize linear SOS for nonnegative projective
quadratics. The map $p\mapsto pp^T$ is quadratic and loses orientation;
these assertions do not themselves give the needed affine-linear section
of the oriented exterior-power mass ball. We make no first
Grassmann-related no-shadow claim. Sanyal--Sottile--Sturmfels
\cite[Theorem 7.6]{orbitopes} exclude an unprojected spectrahedron for
$G_{3,6}$; their discussion distinguishes projected lifts. That weaker
obstruction is not used here.

The verification package includes the complete pinned original proof and
coordinate certificate, independently rederived geometry and moment/source
audits, bounded priority records and exact arithmetic checkers. The
coordinate work checks reality, rank 225, both exposing functionals, the
full polynomial projector identity and mixed Pl\"ucker/minor relations.
Hostile orientation, support, phase and minor cases test fragile boundaries.
Finite computations corroborate the elementary bridge; they do not reprove
Fawzi, certify the universal conclusion or produce an explicit failing form.

The priority audit used 84 focused queries and 25 primary PDF retrievals,
including multiple versions. No inspected body supplied the complete
target assembly or negative conclusion, but unindexed or differently
formulated work remains possible. This is a bounded no-match finding,
not proof of originality or current worldwide openness. The K\"ahler
exposing form itself is a classical positive example, not the existential
failing calibration. All lower-dimensional cases remain unaddressed here.

AI tools were used extensively in solving, verifying, source auditing and
preparation. This preprint is unrefereed; independent adversarial reviews
are AI reviews. Conventional human peer review, journal acceptance and
formal proof-assistant certification are not claimed.

\begin{thebibliography}{99}
\bibitem{sanyal} R. Sanyal, Grassmann orbitopes, Oberwolfach Report 13/2015,
707--709. \url{https://ems.press/content/serial-article-files/46561}.
\bibitem{hl} R. Harvey and H. B. Lawson, Jr., Calibrated geometries,
Acta Math. 148 (1982), 47--157. Question 6.5 and Theorem 6.11.
\bibitem{fawzi} H. Fawzi, The set of separable states has no finite
semidefinite representation except in dimension $3\times2$,
Commun. Math. Phys. 386 (2021), 1319--1335.
\url{https://doi.org/10.1007/s00220-021-04163-2}.
First preprint: arXiv:1905.02575v1 (4 May 2019).
\bibitem{barvinok} A. Barvinok and G. Blekherman, Convex geometry of orbits,
MSRI Publications 52 (2005), 51--78, Example 1.3.
\bibitem{vrana} P. Vrana and P. L\'evay, Special entangled quantum systems
and the Freudenthal construction, arXiv:0902.2269v2 (2009), \S3.1.
\bibitem{fermion} J. Grabowski, M. Ku\'s and G. Marmo, Segre maps and
entanglement for multipartite systems of indistinguishable particles,
J. Phys. A 45 (2012), 105301; arXiv:1111.4812v2.
\bibitem{theta} J. Gouveia, P. A. Parrilo and R. R. Thomas, Theta bodies for
polynomial ideals, SIAM J. Optim. 20 (2010), 2097--2118.
\bibitem{lifts} J. Gouveia, P. A. Parrilo and R. R. Thomas, Lifts of convex
sets and cone factorizations, Math. Oper. Res. 38 (2013), 248--264.
\bibitem{survey} H. Fawzi, J. Gouveia, P. A. Parrilo, J. Saunderson and
R. R. Thomas, Lifting for simplicity: concise descriptions of convex sets,
SIAM Review 64 (2022); arXiv:2002.09788v2.
\bibitem{rostalski} P. Rostalski, SDP relaxations for the Grassmann orbitope,
BIRS slides, 15 February 2010.
\url{https://www.birs.ca/workshops/2010/10w5007/files/rostalski.pdf}.
\bibitem{chapter} J. Gouveia and R. R. Thomas, Spectrahedral approximations
of convex hulls of algebraic sets, Chapter 7 in G. Blekherman, P. A. Parrilo
and R. R. Thomas (eds.), Semidefinite Optimization and Convex Algebraic
Geometry, SIAM (2013), especially p.323.
\bibitem{curvature} R. G. Bettiol, M. Kummer and R. A. E. Mendes, Convex
algebraic geometry of curvature operators, SIAM J. Appl. Algebra Geom.
5 (2021), 200--228; arXiv:1908.03713v1 (2019), v2 (2021).
\bibitem{bsv} G. Blekherman, G. G. Smith and M. Velasco, Sums of squares
and varieties of minimal degree, J. Amer. Math. Soc. 29 (2016), 893--913.
\bibitem{orbitopes} R. Sanyal, F. Sottile and B. Sturmfels, Orbitopes,
Mathematika 57 (2011), 275--314.
\end{thebibliography}
\end{document}
